\documentclass{article}
\usepackage{amsmath,amssymb}
\title{Deriving joint-witness products in the original native presentation}
\date{15 November 2026}
\begin{document}\maketitle
\paragraph{Status.} Independently reviewed independent-input construction. It supplies the local-product hypotheses of the shared-interaction theorem inside the original comma presentation. It does not replace that presentation by a span model or assume its desired native products. The conditions are sufficient, not asserted necessary or automatically true of every O1 instance.

\paragraph{Inputs before the mixed hom.} Fix a pure category $\mathcal C$ with finite products and pullbacks, valid context and content presentations, and their already established mixed interaction. At interfaces $A,B$ write
\[
 e_A:\mathcal M_A\to\mathcal C,\qquad d_B:\mathcal N_B\to\mathcal C,
 \qquad H(A,B)=(e_A\downarrow d_B).
\]
Require finite products in $\mathcal M_A$ and $\mathcal N_B$. Require $e_A$ to admit cartesian lifts of maps into its extension objects. It suffices to admit the specific pullback projections used below, but a chosen fibration $e_A$ is a simpler independent sufficient hypothesis. This concerns the extension functor, not merely the context projection to the pure base. No preservation of products by $e_A$ or $d_B$, and no faithfulness of either, is assumed.

\paragraph{Construction.} For a finite family $(P_i,Q_i,f_i:e_AP_i\to d_BQ_i)$, choose
\[
 P=\prod_iP_i,\qquad Q=\prod_iQ_i,
 \qquad Z=\prod_i d_BQ_i.
\]
The context projections and bodies give $v:e_AP\to Z$. The content projections give the canonical map $\delta:d_BQ\to Z$, not assumed invertible. Form
\[
 T=e_AP\times_Zd_BQ,
 \qquad t:T\to e_AP,\quad b:T\to d_BQ.
\tag{1}
\]
Choose the $e_A$-cartesian lift $c:P^*\to P$ of $t$, identifying $e_AP^*$ with $T$ through its chosen extension comparison. The claimed native product is $(P^*,Q,b)$. Its projections have context maps $P^*\xrightarrow cP\to P_i$ and content maps $Q\to Q_i$. Their native body equations are exactly the pullback square (1). The empty family uses the terminal objects of both grade categories and $Z=1_{\mathcal C}$; the same construction supplies a terminal mixed object.

\paragraph{Full native universal property.} Let $(R,S,g)$ have a cone of actual native arrows $(r_i,s_i)$ to the family. Grade-product universality supplies unique $r:R\to P$ and $s:S\to Q$. The component body equations imply
\[
 v\,e_A r=\delta\,d_Bs\,g.
\]
Therefore the pullback gives a unique $k:e_AR\to T$ with $tk=e_Ar$ and $bk=d_Bs\,g$. Cartesian universality of $c$ gives a unique $r^*:R\to P^*$ above $k$ such that $cr^*=r$. The pair $(r^*,s)$ has the required mixed body equation and projects to the supplied cone.

Conversely every factorization must have those grade maps $r,s$ after projection, that extension map $k$ by pullback uniqueness, and then that context arrow $r^*$ by cartesian uniqueness. Thus it is unique as an actual native arrow, even if distinct context arrows can have equal extension images in general. This proves the product in $H(A,B)$ rather than merely a product of its body observations. Canonical product comparisons and their coherence are determined by these native projections.

\paragraph{Pure input changes: an independent exactness condition.} For $a:A'\to A$, let the original context reindexing $a^*:\mathcal M_A\to\mathcal M_{A'}$ have its actual extension natural transformation
\[
 r_a:e_{A'}a^*\Rightarrow e_A.
\]
Assume $a^*$ preserves context-grade finite products. Assume it preserves the chosen extension-cartesian lifts used in (1), and that the naturality square of $r_a$ at each such lift $c:P^*\to P$ is a pullback:
\[
 \begin{array}{ccc}
 e_{A'}a^*P^*&\longrightarrow&e_AP^*\\
 \downarrow&&\downarrow t\\
 e_{A'}a^*P&\xrightarrow{r_a(P)}&e_AP.
 \end{array}
\tag{2}
\]
These are conditions on context reindexing and its extension, before the mixed hom product. The content grades are unchanged by pure input restriction. Their comparison maps are those of the original context action, not a new target product-preservation axiom.

By context-product preservation, the first grade product in the restricted family is $a^*P$, with its chosen comparison. The bottom map to $Z$ is $v r_a(P)$. Pullback pasting in (1)--(2) identifies the extension of $a^*P^*$ with
$e_{A'}a^*P\times_Zd_BQ$. Preservation of the cartesian lift then identifies its actual context grade with the grade constructed for that restricted family, not just its extension. The bodies agree by their two pullback projections. Hence the original pure-input functor
\[
 H(A,B)\longrightarrow H(A',B)
\]
preserves finite native products through the canonical comparison. Identity and composite input changes satisfy the corresponding equations by the original action coherence and the unique native product maps. No claim of product preservation for arbitrary pure output transport is made.

\paragraph{Consequent interaction, not a postulated mixed law.} The realized process bicategory of the original interaction now has the local finite products just proved and the pure-input preservation required by the shared-interaction construction. To apply the shared-state theorem, require every chosen state-action input arrow $E_A(h)$ to be an original pure map covered by these reindexing/BC hypotheses, or independently verify product preservation for that arrow class. Under this condition its joint-witness native state operation and forward sequential comparison follow from universal projections. No preservation for newly added pure modules or arbitrary state-action arrows in a later completion is inferred. The sequential composition here is the original one induced by its established crossing; no new compatibility cell between crossing and product is postulated. Functoriality of that actual composition acts on the product projections and supplies the forward comparison. This does not rederive the crossing or relax its original hypotheses.

Output transport still has only its canonical forward product comparison. Consequently synchronization of this joint operation remains directed in general, with strong compatibility only under the separately stated output exactness at the actual cone projections. The local product unit is the terminal hom process, not the original chronological identity, and is not silently substituted for the tensor unit.

\paragraph{Recovery of an original family.} For the display context presentation with all maps, take $\mathcal M_A=\mathcal C/A$ and $e_A$ the domain functor. Its product is pullback over $A$, and extension-cartesian lifting is composition of a map into an apex with its map to $A$. Use the trivial content fibre with its single extension object $B$. Then $d_B$ generally does not preserve products: the canonical map is the diagonal $B\to B^m$. Nevertheless (1) constructs the fibre product of span witnesses over $A\times B$, including the terminal span with apex $A\times B$ for the empty family. Context pullback along $a$ satisfies (2) by pullback pasting and preserves fibre products. Thus the theorem recovers the actual native span family already present in the original unifier, with all apex maps retained.

A restricted display class must be closed under the extra restriction maps (1) and the required grade products for this proof to apply. A nonfaithful augmented context family likewise needs its own extension-cartesian lift verification; the theorem does not infer it from being a contextad. Other content fibres need actual grade products. These are explicit structural boundaries, not choices of concrete operational examples.

\paragraph{Attribution and purpose.} Comma-limit and cartesian-lift universal properties are standard. The proof is given directly to avoid a mistaken inference that base products or content products alone suffice. In particular the simple sufficient condition that $d_B$ preserve products would exclude even the trivial content presentation at a general $B$. The extension-lift construction both avoids that restriction and shows exactly which independent context compatibility replaces it. No novelty is claimed for comma limits, spans, or pullback pasting.
\end{document}
